\section{Conclusions}
The 81-paper corpus establishes with reasonable confidence that marine biofouling control is a system problem rather than a single-material problem. Surface chemistry, wettability, roughness, hydration, lubrication, active particles, and mechanical compliance can influence attachment and release, but the direction and magnitude depend on fouling stage, organism, motion, ageing, and cleaning. Nanocomposites and bioinspired interfaces are promising at material and coupon scale; the evidence is weaker for long-duration, full-scale, abrasion-aware deployment.

The engineering consequence is that antifouling, fouling-release, anti-adhesion, self-cleaning, and removal should be selected as distinct control modes. Hydrodynamic and operating-ship studies show why biofilm and roughness matter for drag, powering, fuel, and emissions, while cleaning studies show that maintenance can both recover performance and create coating-wear and wastewater burdens. Thus a biocide-free label alone cannot establish sustainability.

The most defensible pathway is an integrated architecture: select a surface for the asset and operating profile; verify retained adhesion and roughness after immersion and cleaning; monitor fouling condition; apply targeted, low-damage cleaning; contain released material; and evaluate life-cycle cost and environmental exposure. Future research should replace isolated laboratory rankings with standardized, multi-stage validation that follows the chain from interface to fouling, hydrodynamics, maintenance, release, and life-cycle outcome.
